% Standalone resolution manuscript generated from solution.md.
% Compile with XeLaTeX; author and review metadata are maintained in Markdown.
\documentclass[10pt,a4paper]{article}
\usepackage[margin=24mm,headheight=15pt]{geometry}
\usepackage{fontspec}
\setmainfont{DejaVuSerif.ttf}[BoldFont=DejaVuSerif-Bold.ttf,ItalicFont=DejaVuSerif-Italic.ttf,BoldItalicFont=DejaVuSerif-BoldItalic.ttf]
\setsansfont{DejaVuSans.ttf}[BoldFont=DejaVuSans-Bold.ttf,ItalicFont=DejaVuSans-Oblique.ttf]
\setmonofont{DejaVuSansMono.ttf}
\usepackage{amsmath,amssymb,unicode-math}
\setmathfont{latinmodern-math.otf}
\usepackage{xcolor,microtype,parskip,needspace,fancyhdr}
\usepackage{bookmark}
\usepackage{xurl}
\hypersetup{colorlinks=true,urlcolor=blue,linkcolor=blue,pdftitle={SP-13:
Trace-norm-small perturbations preserve Hermitian spectral
distributions},pdfauthor={George Stepaniants},pdflang={en-GB}}
\urlstyle{same}
\setlength{\emergencystretch}{3em}
\providecommand{\tightlist}{\setlength{\itemsep}{0pt}\setlength{\parskip}{0pt}}
\providecommand{\pandocbounded}[1]{#1}
\setcounter{secnumdepth}{0}
\allowdisplaybreaks[1]
\widowpenalty=10000
\clubpenalty=10000
\pagestyle{fancy}
\fancyhf{}
\fancyhead[L]{\small\sffamily RESOLUTION / SP-13}
\fancyhead[R]{\small\sffamily 11 September 2026}
\fancyfoot[L]{\parbox[b]{0.9\textwidth}{\fontsize{8}{10}\selectfont AI-assisted
proof; the exact-source reviews and their limits are documented in the
submission record.}}
\fancyfoot[R]{\footnotesize\thepage}
\begin{document}
{\raggedright\Large\bfseries SP-13: Trace-norm-small perturbations
preserve Hermitian spectral distributions\par}
\vspace{0.4em}
{\normalsize\bfseries George Stepaniants\par}
{\small Department of Computing and Mathematical Sciences, California
Institute of Technology, Pasadena, California, USA.\par}
\vspace{0.6em}
\pagestyle{plain}

This manuscript proves the complete Barbarino-Serra-Capizzano trace-norm
perturbation conjecture. It was developed with substantial ChatGPT/Codex
assistance at the author's request. A separate
\href{https://github.com/ajt60gaibb/OpenProblemsInNLA/blob/main/references/stepaniants-sp13-2026-09-11/verification/SP-13-independent-review.md}{independent
Codex-agent mathematical review} returned PASS for the complete target;
a
\href{https://github.com/ajt60gaibb/OpenProblemsInNLA/blob/main/references/stepaniants-sp13-2026-09-11/verification/SP-13-root-math-review.md}{coordinating-agent
audit} also found no gap. No external human peer review, formal
proof-assistant verification, or priority determination is asserted. The
\href{https://github.com/ajt60gaibb/OpenProblemsInNLA/blob/main/references/stepaniants-sp13-2026-09-11/README.md}{submission
record} preserves the frozen source, original target and verification
evidence.

\subsection{1. Exact target}\label{exact-target}

For \(A_n\in\mathbb C^{n\times n}\) and measurable
\(f:[0,1]\to\mathbb R\), the notation \(\{A_n\}\sim_\lambda f\) means

\[\lim_{n\to\infty}\frac1n\sum_{j=1}^n F(\lambda_j(A_n))
=\int_0^1 F(f(t))\,dt\qquad(F\in C_c(\mathbb C)).\]

The functions are continuous, complex-valued, and compactly supported.
Eigenvalues have algebraic multiplicity. The trace norm is
\(\|X\|_* = \sum_j\sigma_j(X)\); \(X^*\) is the conjugate transpose.

\textbf{Theorem.} Suppose \(H_n=H_n^*\), \(\{H_n\}\sim_\lambda f\), and
\(E_n\in\mathbb C^{n\times n}\) satisfies \(\|E_n\|_*/n\to0\). Then

\[\{H_n+E_n\}\sim_\lambda f.\]

There is no spectral-norm boundedness assumption on either sequence and
no normality assumption on \(H_n+E_n\). This is the complete canonical
SP-13 target and Conjecture 1 in Section 6, printed page 29, of
Barbarino--Serra-Capizzano {[}BSC20{]}.

\subsection{2. The external weak-type
theorem}\label{the-external-weak-type-theorem}

For a finite matrix \(X\) and \(s>0\), set

\[n_X(s)=\#\{j:\sigma_j(X)>s\},\qquad
\|X\|_{1,\infty}=\sup_{s>0}s\,n_X(s).\]

Let \(\mathcal T_+(X)\) retain the strictly upper-triangular entries of
\(X\) and replace all other entries by zero.

\textbf{Published weak-type estimate.} There is an absolute constant
\(C\) such that, for every matrix size and every complex matrix \(X\),

\[n_{\mathcal T_+(X)}(s)\le \frac{C\|X\|_*}{s}\qquad(s>0).\tag{1}\]

This is the finite-matrix specialization of Randrianantoanina {[}R02,
Theorem 4.8, printed page 23{]}. The source defines the weak-\(L^1\)
quasi-norm on printed page 12 by the distribution function used above.
Its theorem applies to an arbitrary finite family of mutually orthogonal
projections and has an \textbf{absolute} constant. Take the algebra
\(M_n(\mathbb C)\) with its ordinary, unnormalized matrix trace, and its
rank-one coordinate projections. The source uses strictly
lower-triangular truncation; reversing the order of the projections
gives \(\mathcal T_+\). Thus (1) is uniform in \(n\) and imposes neither
positivity nor self-adjointness on \(X\).

Theorem 4.8 is explicitly attributed there to Dodds--Dodds--de
Pagter--Sukochev {[}DDPS99, Theorem 1.4{]}, and {[}R02{]} also supplies
a proof from its noncommutative Hilbert-transform estimate. We use this
published result; we do not claim it as a new theorem or replace it by
the weaker logarithmic trace-norm bound.

\subsection{3. Two elementary spectral
facts}\label{two-elementary-spectral-facts}

For arbitrary matrices \(X,Y\) and \(a,b>0\),

\[n_{X+Y}(a+b)\le n_X(a)+n_Y(b).\tag{2}\]

Indeed, remove the singular components of \(X\) larger than \(a\), and
of \(Y\) larger than \(b\). The remaining sum has spectral norm at most
\(a+b\), while the removed sum has rank at most the right side. The
min-max characterization of singular values proves (2). Also

\[n_X(a)\le\frac{\|X\|_*}{a}.\tag{3}\]

\textbf{Lemma.} Let \(A_n,B_n\) be Hermitian matrices of size \(n\). If

\[\frac{n_{A_n-B_n}(s)}n\longrightarrow0\qquad\text{for every fixed }s>0,\tag{4}\]

then for every \(g\in C_c(\mathbb R)\),

\[\frac1n\operatorname{tr}g(A_n)-\frac1n\operatorname{tr}g(B_n)\longrightarrow0.\tag{5}\]

No bound on \(\|A_n\|\) or \(\|B_n\|\) is required.

\textbf{Proof.} Fix \(s>0\). Split the Hermitian matrix \(A_n-B_n\) by
its spectral decomposition into \(R_n+S_n\), where \(R_n\) contains the
eigenvalues with absolute value greater than \(s\). Then

\[\operatorname{rank}R_n=n_{A_n-B_n}(s)=o(n),\qquad\|S_n\|\le s.\]

For Hermitian matrices differing by rank at most \(r\), their eigenvalue
counting functions differ by at most \(r\) at every real threshold. This
follows from the min-max principle by restricting a trial subspace to
the kernel of the rank-\(r\) difference. Consequently, for
\(g\in C_c^1(\mathbb R)\), integration by parts against the finite
eigenvalue counting measures gives

\[\left|\operatorname{tr}g(A_n)-\operatorname{tr}g(B_n+S_n)\right|
\le \operatorname{rank}R_n\int_{\mathbb R}|g'(x)|\,dx.\tag{6}\]

The eigenvalues of \(B_n+S_n\) and \(B_n\), arranged in increasing
order, differ by at most \(s\), again by min-max. With

\[\omega_g(s)=\sup_{|x-y|\le s}|g(x)-g(y)|,\]

their normalized trace difference is at most \(\omega_g(s)\). Thus the
limsup of the absolute value in (5) is at most \(\omega_g(s)\).
Compactly supported continuous functions are uniformly continuous, so
letting \(s\downarrow0\) proves (5) for \(g\in C_c^1(\mathbb R)\). Every
\(g\in C_c(\mathbb R)\) can be uniformly approximated by such functions;
each normalized trace changes by at most the uniform error. This proves
the lemma, also for complex-valued \(g\). \(\square\)

\subsection{4. Schur decomposition and the finite-dimensional
estimates}\label{schur-decomposition-and-the-finite-dimensional-estimates}

Fix \(n\), write \(H=H_n\), \(E=E_n\), and let \(q=\|E\|_*\). Take any
unitary Schur decomposition

\[U^*(H+E)U=T=D+iJ+N,\]

where \(D\) and \(J\) are real diagonal matrices and \(N\) is strictly
upper triangular. The diagonal entries of \(D+iJ\) are exactly the
eigenvalues of \(H+E\), counted with algebraic multiplicity. Set

\[G=U^*HU,\qquad Z=U^*EU,\qquad
K=\operatorname{Im}T=\frac{T-T^*}{2i}.\]

Because \(G\) is Hermitian,

\[K=\operatorname{Im}Z,\qquad \|K\|_*\le\|Z\|_*=q.\tag{7}\]

For a strictly upper-triangular entry, \(K_{ij}=N_{ij}/(2i)\). Therefore

\[N=2i\,\mathcal T_+(K).\]

The dimension-independent estimate (1) yields

\[n_N(s)\le\frac{2Cq}{s}.\tag{8}\]

Applying (2) to \((N+N^*)/2\) gives

\[n_{\operatorname{Re}N}(s)\le2n_N(s)\le\frac{4Cq}{s}.\tag{9}\]

Taking Hermitian parts in \(G+Z=D+iJ+N\) gives

\[D-G=\operatorname{Re}Z-\operatorname{Re}N.\]

Since \(\|\operatorname{Re}Z\|_*\le q\), equations (2), (3), and (9)
show, for every \(s>0\),

\[n_{D-G}(s)
\le n_{\operatorname{Re}Z}(s/2)+n_{\operatorname{Re}N}(s/2)
\le\frac{(2+8C)q}{s}.\tag{10}\]

Finally, the diagonal of \(K\) is \(J\), and

\[\sum_{j=1}^n|J_{jj}|\le\|K\|_*\le q.\tag{11}\]

For clarity, the first inequality needs no bound on a triangular
projection: if \(S\) is the diagonal matrix with entries
\(\operatorname{sgn}(K_{jj})\), then \(\|S\|\le1\) and
\(\sum_j|K_{jj}|=\operatorname{tr}(SK)\le\|K\|_*\). Equivalently it is
trace-norm contractivity of diagonal pinching.

All estimates (7)--(11) hold for every finite \(n\), independently of
\(H\), the Schur ordering, diagonalizability, or the spectral norm of
\(E\).

\subsection{5. Passage to the prescribed test
functions}\label{passage-to-the-prescribed-test-functions}

Return to the sequences, so \(q_n/n\to0\). By (10), for every fixed
\(s>0\),

\[\frac{n_{D_n-U_n^*H_nU_n}(s)}n\longrightarrow0.\]

The lemma implies that \(D_n\) and \(H_n\) have asymptotically equal
normalized traces for all \(g\in C_c(\mathbb R)\). In particular, for a
prescribed \(F\in C_c(\mathbb C)\), the restriction \(g(x)=F(x)\) to the
real axis belongs to \(C_c(\mathbb R)\), and

\[\frac1n\sum_jF((D_n)_{jj})-\frac1n\sum_jF(\lambda_j(H_n))\longrightarrow0.\tag{12}\]

Write \(d_{nj}=(D_n)_{jj}\) and \(y_{nj}=(J_n)_{jj}\). Equation (11)
implies

\[\#\{j:|y_{nj}|>s\}\le\frac{q_n}{s}.\]

Every continuous compactly supported function on \(\mathbb C\) is
uniformly continuous on the whole plane. Defining its modulus
\(\omega_F(s)\) using Euclidean distance in \(\mathbb C\), we obtain

\[\begin{aligned}
&\left|\frac1n\sum_jF(d_{nj}+iy_{nj})-\frac1n\sum_jF(d_{nj})\right|\\
&\hspace{2em}\le\omega_F(s)+2\|F\|_\infty\frac{q_n}{ns}.
\end{aligned}\tag{13}\]

For every fixed \(s>0\) the second term tends to zero. Then
\(s\downarrow0\) makes the first term tend to zero. The values
\(d_{nj}+iy_{nj}\) are the eigenvalues of \(H_n+E_n\) with their
algebraic multiplicities. Combining (12), (13), and the assumed spectral
distribution of \(H_n\) proves

\[\lim_{n\to\infty}\frac1n\sum_jF(\lambda_j(H_n+E_n))
=\int_0^1F(f(t))\,dt.\]

This holds for every \(F\in C_c(\mathbb C)\) and is exactly the asserted
conclusion. \(\square\)

\subsection{6. Scope and attribution}\label{scope-and-attribution}

The proof allows arbitrary complex perturbations and arbitrary real
measurable symbols on the canonical domain. Its only smallness
assumption is \(q_n/n\to0\). Large exceptional singular values and
eigenvalues are handled by rank bounds and compactly supported test
functions, so no uniform spectral bound enters implicitly.

The published weak-type theorem retains its attribution to
Randrianantoanina and Dodds, Dodds, de Pagter and Sukochev. Barbarino
and Serra-Capizzano retain credit for the conjecture and preceding
perturbation results. This proof applies the dimension-independent
triangular-truncation estimate to the imaginary Hermitian part of a
Schur form.

\newpage

\subsection{References}\label{references}

\begin{itemize}
\tightlist
\item
  \textbf{{[}BSC20{]}} G. Barbarino and S. Serra-Capizzano,
  \emph{Non-Hermitian perturbations of Hermitian matrix-sequences and
  applications to the spectral analysis of the numerical approximation
  of partial differential equations}, Numerical Linear Algebra with
  Applications 27 (2020), e2286.
  \href{https://doi.org/10.1002/nla.2286}{DOI},
  \href{https://giovannibarbarino.github.io/doc/articles/NHperturbation.pdf}{author
  PDF}. Conjecture 1, Section 6, printed page 29; the
  spectral-distribution convention appears on printed page 2.
\item
  \textbf{{[}R02{]}} N. Randrianantoanina, \emph{Spectral subspaces and
  non-commutative Hilbert transforms}, Colloquium Mathematicum 91
  (2002), 9--27. \href{https://doi.org/10.4064/cm91-1-2}{DOI},
  \href{https://www.impan.pl/shop/en/publication/transaction/download/product/87479}{publisher
  PDF}. Theorem 4.8, printed page 23, and its proof on page 24;
  weak-\(L^1\) definition on printed page 12.
\item
  \textbf{{[}DDPS99{]}} P. G. Dodds, T. K. Dodds, B. de Pagter and F. A.
  Sukochev, \emph{Lipschitz continuity of the absolute value in preduals
  of semifinite factors}, Integral Equations and Operator Theory 34
  (1999), 28--44. Theorem 1.4 is the result cited by {[}R02, Theorem
  4.8{]}; the application above uses the statement and proof actually
  checked in {[}R02{]}.
\item
  \textbf{{[}B18{]}} G. Barbarino, \emph{Conjectures on Perturbations of
  Hermitian Sequences},
  \href{https://arxiv.org/abs/1808.05555v1}{arXiv:1808.05555v1}, 2018.
  Context for equivalent formulations, not an additional assumption.
\item
  \textbf{{[}BG25{]}} G. Barbarino and C. Garoni, \emph{GLT sequences
  and normal matrices}, Electronic Journal of Linear Algebra 41 (2025),
  1--20.
  \href{https://journals.uwyo.edu/index.php/ela/article/download/8929/6949/23285}{Primary
  PDF}. Theorem 3.2, contextual comparison only.
\end{itemize}
\end{document}
